\documentclass[11pt]{article}
\newcommand{\ListaTitulo}{Gabarito --- Lista 14}
% Preâmbulo compartilhado — MATD48, Listas de Exercícios 2026
% Usado por Lista{NN}.tex e Gabarito{NN}.tex via % Preâmbulo compartilhado — MATD48, Listas de Exercícios 2026
% Usado por Lista{NN}.tex e Gabarito{NN}.tex via % Preâmbulo compartilhado — MATD48, Listas de Exercícios 2026
% Usado por Lista{NN}.tex e Gabarito{NN}.tex via \input{preamble}

\usepackage[utf8]{inputenc}
\usepackage[T1]{fontenc}
\usepackage[brazil]{babel}
\usepackage{fullpage}
\usepackage{amsmath,amssymb}
\usepackage{enumitem}
\usepackage{xcolor}
\usepackage{fancyhdr}
\usepackage{titlesec}
\usepackage{listings}
\usepackage{hyperref}
\usepackage{booktabs}
\usepackage{graphicx}

\definecolor{matd48azul}{HTML}{0B4F6C}
\definecolor{matd48verde}{HTML}{2E8B57}

\titleformat{\section}{\normalfont\Large\bfseries\color{matd48azul}}{\thesection}{1em}{}
\titleformat{\subsection}{\normalfont\large\bfseries\color{matd48azul}}{\thesubsection}{1em}{}

\providecommand{\ListaTitulo}{Lista de Exercícios}

\setlength{\headheight}{14pt}
\pagestyle{fancy}
\fancyhf{}
\lhead{\small MATD48 --- Planejamento de Experimentos A}
\rhead{\small \ListaTitulo}
\cfoot{\thepage}
\renewcommand{\headrulewidth}{0.4pt}

\lstset{
  basicstyle=\ttfamily\small,
  breaklines=true,
  frame=single,
  columns=fullflexible,
  backgroundcolor=\color{gray!8},
  commentstyle=\color{matd48verde},
  keywordstyle=\color{matd48azul}\bfseries,
  language=R,
  extendedchars=true,
  literate=%
    {á}{{\'a}}1 {à}{{\`a}}1 {â}{{\^a}}1 {ã}{{\~a}}1
    {é}{{\'e}}1 {ê}{{\^e}}1 {í}{{\'i}}1 {ó}{{\'o}}1
    {ô}{{\^o}}1 {õ}{{\~o}}1 {ú}{{\'u}}1 {ç}{{\c c}}1
    {Á}{{\'A}}1 {À}{{\`A}}1 {Â}{{\^A}}1 {Ã}{{\~A}}1
    {É}{{\'E}}1 {Ê}{{\^E}}1 {Í}{{\'I}}1 {Ó}{{\'O}}1
    {Ô}{{\^O}}1 {Õ}{{\~O}}1 {Ú}{{\'U}}1 {Ç}{{\c C}}1
}

\hypersetup{colorlinks=true, linkcolor=matd48azul, urlcolor=matd48azul, citecolor=matd48azul}

\newcommand{\cabecalho}[2]{%
  \begin{center}
    {\Large\bfseries\color{matd48azul} #1}\\[0.3em]
    {\large #2}\\[0.3em]
    \rule{\linewidth}{0.6pt}
  \end{center}
  \vspace{0.5em}
}

\newenvironment{questao}[1]{\vspace{1em}\noindent\textbf{Questão #1.}\hspace{0.4em}}{\par}
\newenvironment{solucao}{\vspace{0.3em}\noindent\textit{Solução.}\hspace{0.4em}\color{black!85}}{\par\vspace{0.5em}}


\usepackage[utf8]{inputenc}
\usepackage[T1]{fontenc}
\usepackage[brazil]{babel}
\usepackage{fullpage}
\usepackage{amsmath,amssymb}
\usepackage{enumitem}
\usepackage{xcolor}
\usepackage{fancyhdr}
\usepackage{titlesec}
\usepackage{listings}
\usepackage{hyperref}
\usepackage{booktabs}
\usepackage{graphicx}

\definecolor{matd48azul}{HTML}{0B4F6C}
\definecolor{matd48verde}{HTML}{2E8B57}

\titleformat{\section}{\normalfont\Large\bfseries\color{matd48azul}}{\thesection}{1em}{}
\titleformat{\subsection}{\normalfont\large\bfseries\color{matd48azul}}{\thesubsection}{1em}{}

\providecommand{\ListaTitulo}{Lista de Exercícios}

\setlength{\headheight}{14pt}
\pagestyle{fancy}
\fancyhf{}
\lhead{\small MATD48 --- Planejamento de Experimentos A}
\rhead{\small \ListaTitulo}
\cfoot{\thepage}
\renewcommand{\headrulewidth}{0.4pt}

\lstset{
  basicstyle=\ttfamily\small,
  breaklines=true,
  frame=single,
  columns=fullflexible,
  backgroundcolor=\color{gray!8},
  commentstyle=\color{matd48verde},
  keywordstyle=\color{matd48azul}\bfseries,
  language=R,
  extendedchars=true,
  literate=%
    {á}{{\'a}}1 {à}{{\`a}}1 {â}{{\^a}}1 {ã}{{\~a}}1
    {é}{{\'e}}1 {ê}{{\^e}}1 {í}{{\'i}}1 {ó}{{\'o}}1
    {ô}{{\^o}}1 {õ}{{\~o}}1 {ú}{{\'u}}1 {ç}{{\c c}}1
    {Á}{{\'A}}1 {À}{{\`A}}1 {Â}{{\^A}}1 {Ã}{{\~A}}1
    {É}{{\'E}}1 {Ê}{{\^E}}1 {Í}{{\'I}}1 {Ó}{{\'O}}1
    {Ô}{{\^O}}1 {Õ}{{\~O}}1 {Ú}{{\'U}}1 {Ç}{{\c C}}1
}

\hypersetup{colorlinks=true, linkcolor=matd48azul, urlcolor=matd48azul, citecolor=matd48azul}

\newcommand{\cabecalho}[2]{%
  \begin{center}
    {\Large\bfseries\color{matd48azul} #1}\\[0.3em]
    {\large #2}\\[0.3em]
    \rule{\linewidth}{0.6pt}
  \end{center}
  \vspace{0.5em}
}

\newenvironment{questao}[1]{\vspace{1em}\noindent\textbf{Questão #1.}\hspace{0.4em}}{\par}
\newenvironment{solucao}{\vspace{0.3em}\noindent\textit{Solução.}\hspace{0.4em}\color{black!85}}{\par\vspace{0.5em}}


\usepackage[utf8]{inputenc}
\usepackage[T1]{fontenc}
\usepackage[brazil]{babel}
\usepackage{fullpage}
\usepackage{amsmath,amssymb}
\usepackage{enumitem}
\usepackage{xcolor}
\usepackage{fancyhdr}
\usepackage{titlesec}
\usepackage{listings}
\usepackage{hyperref}
\usepackage{booktabs}
\usepackage{graphicx}

\definecolor{matd48azul}{HTML}{0B4F6C}
\definecolor{matd48verde}{HTML}{2E8B57}

\titleformat{\section}{\normalfont\Large\bfseries\color{matd48azul}}{\thesection}{1em}{}
\titleformat{\subsection}{\normalfont\large\bfseries\color{matd48azul}}{\thesubsection}{1em}{}

\providecommand{\ListaTitulo}{Lista de Exercícios}

\setlength{\headheight}{14pt}
\pagestyle{fancy}
\fancyhf{}
\lhead{\small MATD48 --- Planejamento de Experimentos A}
\rhead{\small \ListaTitulo}
\cfoot{\thepage}
\renewcommand{\headrulewidth}{0.4pt}

\lstset{
  basicstyle=\ttfamily\small,
  breaklines=true,
  frame=single,
  columns=fullflexible,
  backgroundcolor=\color{gray!8},
  commentstyle=\color{matd48verde},
  keywordstyle=\color{matd48azul}\bfseries,
  language=R,
  extendedchars=true,
  literate=%
    {á}{{\'a}}1 {à}{{\`a}}1 {â}{{\^a}}1 {ã}{{\~a}}1
    {é}{{\'e}}1 {ê}{{\^e}}1 {í}{{\'i}}1 {ó}{{\'o}}1
    {ô}{{\^o}}1 {õ}{{\~o}}1 {ú}{{\'u}}1 {ç}{{\c c}}1
    {Á}{{\'A}}1 {À}{{\`A}}1 {Â}{{\^A}}1 {Ã}{{\~A}}1
    {É}{{\'E}}1 {Ê}{{\^E}}1 {Í}{{\'I}}1 {Ó}{{\'O}}1
    {Ô}{{\^O}}1 {Õ}{{\~O}}1 {Ú}{{\'U}}1 {Ç}{{\c C}}1
}

\hypersetup{colorlinks=true, linkcolor=matd48azul, urlcolor=matd48azul, citecolor=matd48azul}

\newcommand{\cabecalho}[2]{%
  \begin{center}
    {\Large\bfseries\color{matd48azul} #1}\\[0.3em]
    {\large #2}\\[0.3em]
    \rule{\linewidth}{0.6pt}
  \end{center}
  \vspace{0.5em}
}

\newenvironment{questao}[1]{\vspace{1em}\noindent\textbf{Questão #1.}\hspace{0.4em}}{\par}
\newenvironment{solucao}{\vspace{0.3em}\noindent\textit{Solução.}\hspace{0.4em}\color{black!85}}{\par\vspace{0.5em}}


\begin{document}

\cabecalho{Gabarito --- Lista de Exercícios 14}{Fatoriais $3^k$, fracionados $2^{k-p}$ e superfície de resposta (Capítulo 6 / Aula 14)}

\begin{questao}{1} \textbf{(Engenharia --- estrutura de graus de liberdade de um $3^2$)}
\end{questao}
\begin{solucao}
\begin{enumerate}[label=(\alph*)]
  \item Cada fator de 3 níveis tem $3-1=\mathbf{2}$ graus de liberdade. A interação
    $A{\times}B$ tem $(3-1)(3-1)=2\times2=\mathbf{4}$ graus de liberdade (produto dos graus de
    liberdade dos dois fatores, mesma regra do fatorial A×B geral do Capítulo 5).
  \item
  \begin{center}
  \begin{tabular}{lc}
  \toprule
  Fonte & gl \\ \midrule
  Velocidade & 2 \\
  Ângulo & 2 \\
  Velocidade$\times$Ângulo & 4 \\
  Erro & $ab(r-1) = 9(3) = 27$ \\
  Total & $abr-1 = 35$ \\ \bottomrule
  \end{tabular}
  \end{center}
  \item Um $2^2$ estimaria apenas a tendência \textbf{linear} de cada fator (e a interação
    linear$\times$linear), porque com só 2 níveis por fator não há informação suficiente para
    detectar \textbf{curvatura} (efeito quadrático). O $3^2$, com 3 níveis por fator, permite
    decompor os 2 graus de liberdade de cada efeito principal em componente linear e componente
    quadrático (polinômios ortogonais, Capítulo 3) --- e é exatamente essa curvatura que a
    metodologia de superfície de resposta (Seção 6.6) explora.
\end{enumerate}
\end{solucao}

\begin{questao}{2} \textbf{(Dedução --- estrutura de aliases de um fracionado $2^{4-1}$)}
\end{questao}
\begin{solucao}
\begin{enumerate}[label=(\alph*)]
  \item Multiplicando $I=ABCD$ por $A$: $A\cdot I = A\cdot ABCD \Rightarrow A = A^2BCD = BCD$
    (pois $A^2=I$ e $I$ multiplicando qualquer efeito não o altera). Logo \textbf{$A$ é aliasado
    com $BCD$}.
  \item Multiplicando por $AB$: $AB\cdot I = AB\cdot ABCD \Rightarrow AB = A^2B^2CD = CD$. Logo
    \textbf{$AB$ é aliasado com $CD$}.
  \item Multiplicando por $C$: $C\cdot I = C\cdot ABCD \Rightarrow C = ABC^2D = ABD$. Logo
    \textbf{$C$ é aliasado com $ABD$}.
  \item A relação definidora $I=ABCD$ tem 4 letras, então este é um fracionado de
    \textbf{resolução IV}. Isso garante que todo efeito principal fica aliasado apenas com
    interações de \emph{três} fatores (item a: $A$ com $BCD$) --- nunca com interações duplas ---
    de modo que, sob a suposição usual de que interações triplas são desprezíveis, os efeitos
    principais são estimados "limpos". O preço é que as interações duplas ficam aliasadas
    \emph{entre si} (item b: $AB$ com $CD$), como esperado de uma resolução IV (nem tão ruim
    quanto III, nem tão boa quanto V).
\end{enumerate}
\end{solucao}

\begin{questao}{3} \textbf{(Ciência de dados/Engenharia --- ponto estacionário de uma superfície de resposta)}
\end{questao}
\begin{solucao}
\begin{enumerate}[label=(\alph*)]
  \item
  $$\frac{\partial \hat y}{\partial x_1} = 4 - 4x_1 + 2x_2, \qquad
  \frac{\partial \hat y}{\partial x_2} = -6 - 6x_2 + 2x_1.$$
  \item Igualando a zero: $4-4x_1+2x_2=0 \Rightarrow 2x_1-x_2=2$; e $-6-6x_2+2x_1=0 \Rightarrow
    x_1-3x_2=3$. Da primeira, $x_2=2x_1-2$. Substituindo na segunda:
    $x_1-3(2x_1-2)=3 \Rightarrow x_1-6x_1+6=3 \Rightarrow -5x_1=-3 \Rightarrow x_1^*=0{,}6$, e
    $x_2^*=2(0{,}6)-2=-0{,}8$. Ponto estacionário: $\boldsymbol{(0{,}6,\,-0{,}8)}$.
  \item
  $$\mathbf{H} = \begin{bmatrix} \partial^2\hat y/\partial x_1^2 &
  \partial^2\hat y/\partial x_1\partial x_2 \\
  \partial^2\hat y/\partial x_1\partial x_2 & \partial^2\hat y/\partial x_2^2 \end{bmatrix} =
  \begin{bmatrix} -4 & 2 \\ 2 & -6 \end{bmatrix}.$$
  Autovalores: $\det(\mathbf{H}-\lambda\mathbf{I})=(-4-\lambda)(-6-\lambda)-4 =
  \lambda^2+10\lambda+20=0 \Rightarrow \lambda = \dfrac{-10\pm\sqrt{100-80}}{2} =
  -5\pm\sqrt5 \approx \mathbf{-2{,}76 \text{ e } -7{,}24}$.
  \item Os dois autovalores são \textbf{negativos} $\Rightarrow$ Hessiana negativa definida
    $\Rightarrow$ o ponto estacionário é um \textbf{máximo}. Substituindo $(0{,}6,-0{,}8)$ no
    modelo:
    $$\hat y = 80+4(0{,}6)-6(-0{,}8)-2(0{,}6)^2-3(-0{,}8)^2+2(0{,}6)(-0{,}8) = 80+2{,}4+4{,}8-0{,}72-1{,}92-0{,}96 = \mathbf{83{,}6}.$$
\end{enumerate}
\end{solucao}

\begin{questao}{4} \textbf{(Confusão em $3^k$ --- comparando os geradores $AB$ e $AB^2$)}
\end{questao}
\begin{solucao}
\begin{enumerate}[label=(\alph*)]
  \item
  \begin{center}
  \begin{tabular}{c|ccccccccc}
  $(A,B)$ & (0,0) & (1,0) & (2,0) & (0,1) & (1,1) & (2,1) & (0,2) & (1,2) & (2,2) \\ \hline
  $L_{AB^2}=(A+2B)\bmod3$ & 0 & 1 & 2 & 2 & 0 & 1 & 1 & 2 & 0
  \end{tabular}
  \end{center}
  Blocos sob $AB^2$: Bloco 0 $=\{(0,0),(1,1),(2,2)\}$; Bloco 1 $=\{(1,0),(2,1),(0,2)\}$;
  Bloco 2 $=\{(2,0),(0,1),(1,2)\}$.
  \item
  \begin{center}
  \begin{tabular}{c|ccccccccc}
  $(A,B)$ & (0,0) & (1,0) & (2,0) & (0,1) & (1,1) & (2,1) & (0,2) & (1,2) & (2,2) \\ \hline
  $L_{AB}=(A+B)\bmod3$ & 0 & 1 & 2 & 1 & 2 & 0 & 2 & 0 & 1
  \end{tabular}
  \end{center}
  Blocos sob $AB$: Bloco 0 $=\{(0,0),(2,1),(1,2)\}$; Bloco 1 $=\{(1,0),(0,1),(2,2)\}$;
  Bloco 2 $=\{(2,0),(1,1),(0,2)\}$.
  \item Sob $AB^2$, $(0,0)$ e $(1,1)$ estão no \textbf{mesmo} bloco (Bloco 0, junto com $(2,2)$).
    Sob $AB$, $(0,0)$ está no Bloco 0 e $(1,1)$ está no Bloco 2 --- blocos
    \textbf{diferentes}. Isso mostra concretamente que $AB$ e $AB^2$ são contrastes de grau único
    genuinamente \textbf{distintos}: cada um particiona as 9 combinações de um jeito diferente, e
    escolher um ou outro como gerador de confusão sacrifica informações diferentes --- não são
    formas equivalentes de escrever a "mesma" interação.
\end{enumerate}
\end{solucao}

\begin{questao}{5} \textbf{(Ciência de dados/Engenharia --- exercício computacional em R)}
\end{questao}
\begin{solucao}
\begin{enumerate}[label=(\alph*)]
  \item
  \begin{lstlisting}
modelo <- lm(energia ~ Velocidad + angulo + I(Velocidad^2) + I(angulo^2) +
               Velocidad:angulo, data = energia)

grade <- expand.grid(
  Velocidad = seq(min(energia$Velocidad), max(energia$Velocidad), length.out = 100),
  angulo    = seq(min(energia$angulo), max(energia$angulo), length.out = 100)
)
grade$pred <- predict(modelo, newdata = grade)

grade[which.min(grade$pred), ]
  \end{lstlisting}
  \item Como visto no Capítulo 6 (Seção "Introdução à superfície de resposta"), o ponto
    estacionário deste modelo específico é uma \textbf{sela}: os autovalores da Hessiana têm
    sinais opostos. Uma sela nunca é máximo nem mínimo --- a superfície sobe em uma direção e
    desce em outra a partir dela --- por isso é absolutamente consistente que o mínimo
    \emph{observado na grade} não coincida com o ponto estacionário: o mínimo real, quando o
    ponto estacionário é uma sela, fica necessariamente na \textbf{fronteira} da região
    experimental, não no interior.
\end{enumerate}
\end{solucao}

\end{document}
